\documentclass[aspectratio=169,11pt]{beamer}

\usepackage[T1]{fontenc}
\usepackage[utf8]{inputenc}
\usepackage[french]{babel}
\usepackage{lmodern}
\usepackage{booktabs}
\usepackage{tabularx}

\definecolor{projectblue}{HTML}{17365D}
\definecolor{projectteal}{HTML}{087E8B}
\usetheme{Madrid}
\usecolortheme{default}
\setbeamercolor{structure}{fg=projectblue}
\setbeamercolor{frametitle}{bg=projectblue,fg=white}
\setbeamercolor{block title}{bg=projectteal,fg=white}
\setbeamercolor{block body}{bg=projectteal!7}
\setbeamertemplate{navigation symbols}{}
\setbeamertemplate{footline}{%
  \leavevmode\hbox{%
    \begin{beamercolorbox}[wd=.82\paperwidth,ht=2.7ex,dp=1.2ex,leftskip=1em]{author in head/foot}%
      VAP SSR 2026--2027 \hfill Réunion de lancement -- 6 octobre 2026
    \end{beamercolorbox}%
    \begin{beamercolorbox}[wd=.18\paperwidth,ht=2.7ex,dp=1.2ex,center]{date in head/foot}%
      \insertframenumber/\inserttotalframenumber
    \end{beamercolorbox}}%
}
\setbeamertemplate{itemize items}[circle]

\title[Démonstrateur ML-NIDS]{Démonstrateur d'attaques contre un ML-NIDS}
\subtitle{Réunion de début de projet}
\author{Projet VAP SSR 2026--2027}
\institute{Encadrant : Gregory Blanc}
\date{6 octobre 2026}

\begin{document}

\begin{frame}
  \titlepage
\end{frame}

\begin{frame}{Notre compréhension du projet}
  \begin{block}{Objectif}
    Un démonstrateur pédagogique pour comprendre les attaques contre un ML-NIDS
    et observer leurs effets.
  \end{block}
  \textbf{Livrables attendus}
  \begin{itemize}
    \setlength{\itemsep}{.8em}
    \item Bref état de l'art des attaques contre les ML-NIDS.
    \item Logiciel de démonstration en trois parties :
      \begin{enumerate}
        \setlength{\itemsep}{.5em}
        \item Choix de l'environnement réseau, façon Packet Tracer.
        \item Choix du ML-NIDS.
        \item Choix de l'attaque effectuée.
      \end{enumerate}
  \end{itemize}
  \vfill
  \begin{block}{Une vue pour comprendre l'attaque}
    Visualiser son déroulement en direct ou avec une comparaison avant/après.
  \end{block}
\end{frame}

\begin{frame}[plain]
  \centering
  {\fontsize{54}{64}\selectfont\bfseries\color{projectblue} JIRA\par}
\end{frame}

\begin{frame}{ML-NIDS et attaques concrètes identifiés}
  \small
  \renewcommand{\arraystretch}{1.35}
  \begin{tabularx}{\textwidth}{@{}p{.24\textwidth}Xp{.26\textwidth}@{}}
    \toprule
    \textbf{Modèles ciblés} & \textbf{Manipulation adversariale} & \textbf{Papiers sources} \\
    \midrule
    KitNET / autoencodeurs, Isolation Forest
      & Modifier les délais ; ajouter des paquets.
      & Han (2021), Sharon (2022). \\
    MLP
      & Modifier le comportement TCP à partir des caractéristiques importantes.
      & Okada (2025). \\
    MLP, CNN, ResCNN, LSTM, attention
      & Perturbations transférables ou guidées par des requêtes au détecteur.
      & Mao (2025), Zhang (2022). \\
    E-GraphSAGE, LineGraphSAGE
      & Ajouter des communications pour modifier le graphe observé.
      & Venturi (2024). \\
    SVM, Random Forest, MLP
      & Injecter des données mal étiquetées ou modifier les labels d'entraînement.
      & Venkatesan (2021), Das (2025). \\
    \bottomrule
  \end{tabularx}
  \vfill
  \small Ces couples modèle--attaque sont des pistes bibliographiques : leur intégration
  dans le démonstrateur reste à étudier.
\end{frame}

\begin{frame}{Changer le rythme du trafic pour passer inaperçu}
  \textbf{Modèles concernés :} KitNET et d'autres détecteurs d'anomalies.
  \vspace{1em}
  \begin{block}{Principe}
    Le détecteur observe la façon dont les paquets circulent.
    L'attaquant change ce rythme pour que son trafic ressemble davantage
    à une activité normale.
  \end{block}
  \begin{itemize}
    \setlength{\itemsep}{.8em}
    \item \textbf{TrafficManipulator :} espacer les envois et ajouter des paquets.
    \item \textbf{TANTRA :} modifier les délais entre les paquets sans changer leur contenu.
  \end{itemize}
  \vfill
  \textbf{Idée de visualisation :} une chronologie des paquets avant / après,
  accompagnée du score de détection.
  \vfill
  \footnotesize Papiers : \href{https://doi.org/10.1109/JSAC.2021.3087242}{Han et al. (2021), TrafficManipulator}
  ; \href{https://doi.org/10.1109/TIFS.2022.3201377}{Sharon et al. (2022), TANTRA}.
\end{frame}

\begin{frame}{Repérer ce que regarde le détecteur, puis l'adapter}
  \textbf{Modèle concerné :} un réseau de neurones classique (MLP).
  \vspace{1em}
  \begin{block}{Principe}
    Chercher quels aspects d'une communication influencent la décision du modèle,
    puis les modifier pour rendre l'attaque moins visible.
  \end{block}
  \begin{itemize}
    \setlength{\itemsep}{.8em}
    \item Exemple étudié : une attaque qui essaie de nombreux mots de passe.
    \item Changer la manière de se connecter, par exemple en ouvrant
      une connexion différente à chaque tentative.
    \item L'objectif reste le même, mais la communication vue par le détecteur change.
  \end{itemize}
  \vfill
  \textbf{Idée de visualisation :} les connexions et les alertes avant / après.
  \vfill
  \footnotesize Papier : \href{https://doi.org/10.1007/s10207-025-01016-0}{Okada et al. (2025), Xai-driven black-box adversarial attacks on network intrusion detectors}.
\end{frame}

\begin{frame}{Chercher une modification qui trompe un réseau de neurones}
  \textbf{Modèles concernés :} plusieurs types de réseaux de neurones.
  \vspace{1em}
  \begin{block}{Deux façons de chercher une évasion}
    \begin{itemize}
      \setlength{\itemsep}{.8em}
      \item \textbf{Sur un autre modèle :} préparer l'attaque sur un modèle accessible,
        puis essayer de tromper le détecteur cible avec la même modification.
      \item \textbf{Par essais successifs :} tester des modifications et utiliser les
        réponses du détecteur pour ajuster l'attaque.
    \end{itemize}
  \end{block}
  \vfill
  \textbf{Idée de visualisation :} suivre les essais et voir à quel moment
  une activité malveillante est classée comme normale.
  \vfill
  \footnotesize Papiers : \href{https://doi.org/10.1016/j.jnca.2025.104255}{Mao et al. (2025), Exploring transferable adversarial attacks\ldots}
  \\ \href{https://doi.org/10.1109/TNET.2021.3137084}{Zhang et al. (2022), Adversarial Attacks Against Deep Learning-Based NIDS\ldots}
\end{frame}

\begin{frame}{Ajouter des communications pour brouiller le contexte}
  \textbf{Modèles concernés :} réseaux de neurones sur graphes.
  \vspace{1em}
  \begin{block}{Principe}
    Ces modèles regardent aussi qui communique avec qui.
    L'attaquant ajoute des échanges normaux autour d'une machine compromise
    pour modifier le contexte observé par le détecteur.
  \end{block}
  \begin{itemize}
    \setlength{\itemsep}{.8em}
    \item La communication malveillante reste présente.
    \item De nouvelles communications modifient la représentation du réseau.
    \item Le modèle peut alors moins bien reconnaître l'activité malveillante.
  \end{itemize}
  \vfill
  \textbf{Idée de visualisation :} les machines et leurs liens avant / après
  l'ajout des communications.
  \vfill
  \footnotesize Papier : \href{https://arxiv.org/abs/2403.11830}{Venturi et al. (2024), Problem space structural adversarial attacks for NIDS based on Graph Neural Networks}.
\end{frame}

\begin{frame}{Faire apprendre de mauvais exemples au détecteur}
  \textbf{Modèles concernés :} SVM, forêts d'arbres et réseaux de neurones classiques.
  \vspace{1em}
  \begin{block}{Principe : l'empoisonnement}
    L'attaquant fausse une partie des exemples utilisés pour apprendre.
    Le détecteur risque alors d'apprendre de mauvaises règles de décision.
  \end{block}
  \begin{itemize}
    \setlength{\itemsep}{.8em}
    \item Ajouter des exemples trompeurs dans les données d'entraînement.
    \item Ou changer leur étiquette « normal » / « malveillant ».
    \item L'effet apparaît après l'entraînement : il faut pouvoir agir sur ses données.
  \end{itemize}
  \vfill
  \textbf{Idée de visualisation :} comparer les erreurs du détecteur entraîné
  avec des données propres puis avec des données empoisonnées.
  \vfill
  \footnotesize Papiers : \href{https://doi.org/10.1109/MILCOM52596.2021.9652916}{Venkatesan et al. (2021), Poisoning Attacks and Data Sanitization Mitigations\ldots}
  \\ \href{https://doi.org/10.1109/TNSE.2024.3492032}{Das et al. (2025), Poisoning the Well\ldots}
\end{frame}

\begin{frame}{Récapitulatif : ce qu'il reste à approfondir}
  \begin{block}{Étudier les attaques trouvées}
    Comprendre leur fonctionnement précis et leurs conditions d'application.
  \end{block}
  \begin{itemize}
    \setlength{\itemsep}{.8em}
    \item \textbf{Vérifier la faisabilité :} quels modèles, données et codes
      peut-on utiliser dans notre démonstrateur ?
    \item \textbf{Tester les scénarios :} l'attaque trompe-t-elle le détecteur
      tout en conservant son effet sur le réseau ?
    \item \textbf{Étoffer la sélection :} chercher d'autres attaques pour varier
      les mécanismes et les modèles présentés.
    \item \textbf{Choisir les visualisations :} rendre le principe et le résultat
      de chaque attaque compréhensibles.
  \end{itemize}
\end{frame}

\begin{frame}{Suite du projet : les jalons prévus}
  \begin{block}{Jalon actuel}
    État de l'art des attaques contre les ML-NIDS : organiser les références
    et rédiger une synthèse comparative.
  \end{block}
  \begin{center}
    {\Large $\downarrow$}
  \end{center}
  \begin{columns}[T,onlytextwidth]
    \begin{column}{.48\textwidth}
      \begin{block}{Clarification de l'utilisation de Concap}
        Préciser son rôle et son utilisation dans le projet.
      \end{block}
    \end{column}
    \begin{column}{.48\textwidth}
      \begin{block}{En parallèle : maquette de l'architecture logicielle}
        Représenter les trois parties du logiciel et leurs interactions.
      \end{block}
    \end{column}
  \end{columns}
\end{frame}

\begin{frame}{Discussion avec l'encadrant}
  \begin{enumerate}
    \setlength{\itemsep}{1em}
    \item Le corpus réuni couvre-t-il les principales attentes pour l'état de l'art ?
    \item Quels articles ou quelles familles d'attaques approfondir en priorité ?
    \item Quel niveau de détail et quel format attendre pour le bref état de l'art ?
    %\item Quels points éclaircir sur l'utilisation de Concap ?
    %\item Quel contenu attendre pour la première maquette de l'architecture logicielle ?
  \end{enumerate}
\end{frame}

\end{document}
